\section{How Cartels Organize: Communication, Facilitating Practices, and Transfers}\label{sec:12-cartels}

The models of \cref{sec:12-perfect,sec:12-demand,sec:12-public,sec:12-private-costs} treat collusion as an
equilibrium of a non-cooperative repeated game. Firms in those models never meet, never talk and never sign
anything. The cartels that antitrust authorities prosecute look different. Their members meet regularly, agree on
target prices, divide sales, exchange sales figures, hire auditors and pay one another compensation. This section
asks what these practices accomplish, and why firms take the legal risk of using them when tacit collusion is
available in theory.

The answer turns on information. Suppose that monitoring is perfect and that demand and costs are common
knowledge. An agreement can then do little more than select one of the many equilibria of the repeated game.
Once prices, sales or costs are private, communication and transfers change what the cartel can enforce.
\Cref{sec:12-anatomy} describes the common design of explicit cartels and maps each design element to the problem
it solves. \Cref{sec:12-facilitating} studies business practices and market features that relax the incentive
constraints of \cref{sec:12-perfect} without an explicit agreement, including multimarket contact and the threat of
entry. \Cref{sec:12-hs} develops in depth the model
of \citet{HarringtonSkrzypacz2011}. It shows that quotas, self-reported sales and transfers form an equilibrium
when prices and sales are private. \Cref{sec:12-bm-preview} previews a model in which cartels tolerate deliberate
cheating, which \cref{sec:13-imperfect} develops in full.

\subsection{The anatomy of explicit cartels}\label{sec:12-anatomy}

A cartel, as defined in \cref{sec:12-forms}, colludes explicitly: its members coordinate through direct
communication about prices, quantities, customers or bids. The repeated-game models of this chapter apply to
explicit and tacit collusion alike. They describe which outcomes the threat of future competition can enforce,
whatever the means by which firms coordinate. The records of prosecuted cartels raise a second question, which
concerns the organization of the agreement itself.

\citet[pp.~105--108]{MarshallMarx2012} build their account of cartel organization on \citet{Stigler1964}. For them
the central problem of a cartel is to avert secret deviations by its members. A cartel solves it with three
complementary \emph{collusive structures}. \emph{Pricing structures} implement price increases and output
reductions. \emph{Allocation structures} divide the collusive gain, through market shares, territories or customers,
and set up mechanisms to redistribute gains among members. \emph{Enforcement structures} monitor compliance and
establish the threat of punishment. Marshall and Marx stress that each structure is insufficient alone. An agreement
on price without an agreement on allocation and enforcement will almost surely collapse through secret deviations.
A cartel that has suppressed rivalry among its members can then turn to the forces outside it, such as entrants,
producers of substitutes, buyers and suppliers \citep[pp.~143--144]{MarshallMarx2012}. We first state the two
constraints that every cartel must satisfy. We then take four internal elements in turn, treating monitoring
separately from punishment, and attach each to the theory of this chapter. \Cref{tab:12-cartel-elements} summarizes
the mapping, and the subsection ends with the plus factors that these elements leave behind.

\casehead{Two constraints}
\citet[pp.~187--188]{AskerNocke2021} organize the economics of cartels around two constraints. Let $\mathcal C$ be a
set of potential members. Firm $i \in \mathcal C$ joins only if it satisfies the \emph{participation constraint}, also
called the internal stability condition,
\begin{equation}\label{eq:12-participation}
V_i^{\mathrm{comp}} \;\le\; V_i^{\mathrm{coll}} - \eta_i .
\end{equation}
Here $V_i^{\mathrm{coll}}$ is the firm's value inside the cartel. $V_i^{\mathrm{comp}}$ is its value if it stays out
while the other firms in $\mathcal C$ join. The one-time cost $\eta_i$ of joining captures the frictions of organizing.
No court enforces a cartel agreement, so each member must also satisfy the \emph{no-cheating constraint}
\begin{equation}\label{eq:12-no-cheating}
V_i^{\mathrm{coll}} \;\ge\; \pi_i^{\mathrm{dev}} + \delta\,V_i^{\mathrm{ret}} .
\end{equation}
Here $\pi_i^{\mathrm{dev}}$ is the profit in the period of cheating and $V_i^{\mathrm{ret}}$ the continuation value
after retaliation. Constraint \eqref{eq:12-no-cheating} is the incentive constraint of every model so far.
Constraint \eqref{eq:12-participation} is new. It asks whether each firm wants to be inside at all.

The participation constraint can fail even when collusion is very profitable. \citet[pp.~192--193]{AskerNocke2021}
give an example. Six firms with marginal cost $c$ compete in quantities against inverse demand $P(Q) = a - bQ$. Let
$A = \frac{(a - c)^2}{b}$. In the Cournot equilibrium with $M$ symmetric players each earns $\frac{A}{(M + 1)^2}$, so
each of the six firms earns $\frac{A}{49}$. A cartel of all six earns the monopoly profit $\frac{A}{4}$, or $\frac{A}{24}$
per member. Suppose one firm refuses to join. The other five then act as one player in a Cournot duopoly with the
outsider, and each of the two players earns $\frac{A}{9}$. Each insider gets $\frac{A}{45} > \frac{A}{49}$, so the
five still gain by colluding. The outsider gets $\frac{A}{9}$, far more than $\frac{A}{24}$. Every firm would like to be
the one that stays out. Negotiations over who is in and who gets what are therefore hard, and a threat to stay out
unless paid extra is credible. The example mirrors the merger paradox of quantity competition. \citet{BosHarrington2010}
endogenize cartel formation among firms with different capacities. Smaller firms are more likely to stay out. The
cartel sets a price that serves as an umbrella, and the noncartel firms price below it and produce at capacity.

\casehead{Agreeing on price}
Repeated-game theory predicts a whole set of sustainable prices. In the Bertrand supergame every price between
marginal cost and the monopoly price is sustainable once $\delta \ge 1 - \frac{1}{N}$
(\cref{cor:12-bertrand-any-price}). Firms that collude tacitly must somehow land on the same element of this set.
An explicit agreement solves this selection problem directly. The members of the lysine cartel, for example, met
in July 1992. They agreed to raise the price to \$1.05 per pound by September 30, 1992, and to \$1.20 per pound by
December 30, 1992 \citep[p.~2426]{HarringtonSkrzypacz2011}. A notorious invitation to collude shows how a single
proposal can do the same. In 1982 Robert Crandall, chief executive of American Airlines, telephoned Howard Putnam of
Braniff, a rival on routes out of Dallas-Fort Worth. Crandall proposed that Braniff raise its fares by twenty percent
and promised to raise his own the next morning. Putnam recorded the call and turned it over to the Department of
Justice \citep[pp.~188--189]{AskerNocke2021}. The proposal named the price point, a simple rule of thumb across many
fares, and left the threat of continued fare wars implicit. \citet{GreenPorter1984} made a similar point about their
own trigger strategies. In view of their complexity and the close coordination they need, it seemed natural to
assume that the equilibrium arises from an explicit agreement (p.~89, footnote~5).

Tacit coordination is possible, but it can be slow. \citet{ByrneDeRoos2019} study retail gasoline in Perth, Australia,
from 2001 to 2015 \citep[pp.~189--191]{AskerNocke2021}. Stations had to commit to prices a day ahead and post them
on Fuelwatch, a government website meant to help consumers search. Between 2009 and 2013, as Asker and Nocke date the
episode, the BP chain led the market to higher margins and synchronized price jumps. Byrne and de Roos describe a
three-year transition from April 2009 to August 2012 \citep[p.~592]{ByrneDeRoos2019}. BP first raised prices on Wednesdays, and by mid-2010 Caltex, the
next-largest chain, followed with jumps on Thursdays. BP then moved its own jumps to Thursday and occasionally
stopped leading, to test whether the market would coordinate without it. The market settled on weekly price jumps
on Thursdays, followed by daily price cuts of 2 cents per liter between jumps
\citep[pp.~596--597 and 607--608]{ByrneDeRoos2019}. Byrne and de Roos read the episode as coordination learned
through prices, but they state that they cannot rule out explicit communication between firms, such as phone calls
\citep[p.~617]{ByrneDeRoos2019}. \citet[p.~191]{AskerNocke2021} conclude only that market interaction itself may
provide enough communication for coordination.

Laboratory evidence points the same way. \citet{CooperKuhn2014} study a two-stage game that is a reduced form of a
repeated Bertrand game. The first stage is a pricing game with a low, a medium and a high price. The second is a
coordination game with three equilibria, which are Pareto ranked. Without communication, matched subjects reach the
high price in the first stage in fewer than 10 percent of cases. With limited communication the share is at most 50
percent, and it falls below 20 percent as subjects gain experience. In the least restrictive treatment, which allows
unrestricted chat, the share exceeds 80 percent \citep[p.~191]{AskerNocke2021}. Kaplow cites a coordination
experiment with a less regular pattern (\cref{sec:12-kaplow}).

Pricing algorithms raise the question in a new form \citep[p.~192]{AskerNocke2021}. In simulations by
\citet{CalvanoEtAl2020}, agents that learn to play a repeated Bertrand game through a reinforcement-learning
algorithm reach supracompetitive prices. \citet{AssadEtAl2020} find that the inferred adoption of algorithmic
pricing by German gasoline retailers coincided with increases in margins of up to 28 percent. Whether such outcomes
count as agreements is the question of \cref{sec:12-kaplow}.

Agreement is harder when members differ. A low-cost
firm prefers a lower collusive price than a high-cost firm, and asymmetries change who gains most from a deviation
(\cref{sec:12-perfect}). The resulting bargaining problems are frequent in the historical record.
In the case studies surveyed by \citet{LevensteinSuslow2006}, bargaining problems end more cartels than secret
cheating does (\cref{sec:12-success}).

\casehead{Allocating sales}
Most documented cartels agree on an allocation as well as a price. Lysine and citric acid used quotas in
quantities or in market shares (\cref{tab:12-hs-allocations}). The nucleotide cartel of
\cref{sec:13-imperfect-evidence} assigned customers. Other cartels drew exclusive territories, and bidding rings
used bid rotation or knockout auctions (\cref{sec:13-asker}). An allocation does three things. First, it turns the
problem of policing secret price cuts into a problem of policing quantities. A firm that steals customers ends up
with sales above its quota, and sales are easier to count than the discounts behind them. This is the logic of
\citet{Stigler1964}, who saw the detection of secret price cuts through shifts of customers as the central
problem of oligopoly (\cref{sec:12-public}). Second, an allocation fixes the division of the collusive profit,
which removes one source of bargaining. Third, an allocation can raise joint profit by directing production to the
lowest-cost members. \Cref{sec:12-private-costs} shows the limits of this third role. Without communication, firms
with private costs often do best with a rigid price that ignores costs (\cref{prop:12-abs-rigid}). When firms can
report their costs, the cartel can give more of the market to the member with the lowest cost
\citep{AtheyBagwell2001}.

\casehead{Monitoring}
A cartel monitors compliance by gathering information about the prices or sales of its members. In lysine one
member compiled self-reported sales into monthly ``scorecards''. In citric acid outside Swiss auditors checked the
reports. The Sugar Institute of \cref{sec:13-sugar-institute} collected reports through a trade association, and
many cartels relied on published price lists. \citet[p.~106]{MarshallMarx2012} distinguish two kinds of deviation
that monitoring must handle. \emph{Unintentional deviations} arise, for example, when a sales staff that does not
know of the cartel bids too aggressively or wins a customer it should not have won. If management acknowledges such
mistakes quickly, the cartel can correct them by reallocating future contracts, by subcontracting, or by
end-of-period transfers. \emph{Secret deviations} are special deals that the perpetrator does not acknowledge, and
uncovering them requires costly intelligence. Many cartels delegate part of the monitoring to outsiders. Marshall
and Marx review the twenty-two decisions on industrial-product cartels that the European Commission issued from 2000
through 2005 (pp.~125--127). In eleven of them, the decision mentions a third-party facilitator, such as a trade
association or a consulting or accounting firm (their Table~6.1, our count). The Swiss consultancy AC-Treuhand, which
served several European chemical cartels, was itself prosecuted for its services to the heat-stabilizers cartel
\citep[p.~197]{AskerNocke2021}. The information structure decides which theory applies.
In \citet{GreenPorter1984}, firms see a public price that depends on hidden outputs and on demand
(\cref{sec:12-public}). In \citet{AtheyBagwellSanchirico2004}, prices are public and costs are private
(\cref{sec:12-private-costs}). In the industrial cartels studied by \citet{HarringtonSkrzypacz2011}, both prices
and sales were private, so monitoring had to rely on what firms told one another (\cref{sec:12-hs}).

\casehead{Enforcement and transfers}
The basic enforcement device of repeated-game theory is a price war that hurts every member. A war destroys
profit, and it punishes the innocent along with the guilty. Many cartels also used targeted compensation. Under
the lysine cartel's ``guaranteed buy-ins'' and the citric acid cartel's ``buy-back system'', a firm that sold more
than its quota bought output from members below quota. The International Steel Agreement of 1926 ran a common fund
into which over-quota producers paid (\cref{sec:12-hs} gives the details). An interfirm sale at a price above cost transfers surplus from
the buyer to the seller. Such sales also aroused little suspicion, because they were common in these industries
\citep[p.~2433]{HarringtonSkrzypacz2011}. Transfers of this kind are \emph{asymmetric punishments}. They penalize
the member that sold too much and leave industry profit intact. \citet{HarringtonSkrzypacz2007} first argued that
asymmetric punishments are essential to effective collusion when firms observe one another's sales
\citep[p.~2442]{HarringtonSkrzypacz2011}. \Cref{sec:12-hs} extends the argument to sales that are private.

\begin{table}[htbp]
\centering
\caption{The design elements of explicit cartels}
\label{tab:12-cartel-elements}
\small
\begin{threeparttable}
\begin{tabularx}{\textwidth}{@{}>{\raggedright\arraybackslash}p{2.0cm}
  >{\raggedright\arraybackslash}X >{\raggedright\arraybackslash}X >{\raggedright\arraybackslash}p{2.9cm}@{}}
\toprule
Element & Problem it solves & Documented examples & Theory in this book \\
\midrule
Price & Select one price from the set of sustainable prices. Settle the level among members with different
interests. & Lysine targets of \$1.05 and \$1.20 per pound (1992) & \cref{cor:12-bertrand-any-price},
\cref{sec:12-success} \\
\addlinespace
Allocation & Turn secret price cuts into visible excess sales. Divide the gains. Direct output to efficient
members. & Lysine and citric acid quotas (\cref{tab:12-hs-allocations}). Customer allocation (nucleotides). Bid
rotation and knockouts (stamp ring). & \cref{sec:12-public,sec:12-private-costs,sec:12-hs} \\
\addlinespace
Monitoring & Verify compliance when prices or sales are private. & Monthly sales scorecards (lysine). Swiss
auditors (citric acid). Trade association rules and reports (Sugar Institute). &
\cref{sec:12-public,sec:12-hs,sec:13-information} \\
\addlinespace
Enforcement & Deter deviations at the least cost to the cartel. & Guaranteed buy-ins (lysine). Buy-backs (citric
acid). A fund of \$4 and \$2 per ton (steel, 1926). Price wars (railroads, 1880s). &
\cref{sec:12-perfect,sec:12-hs,sec:13-jec} \\
\bottomrule
\end{tabularx}
\begin{tablenotes}\footnotesize
\item Notes: Our synthesis. The examples come from \citet{HarringtonSkrzypacz2011}, \citet{GenesoveMullin2001},
\citet{BernheimMadsen2017} and \citet{Asker2010}, which \cref{ch:collusion-evidence} discusses.
\end{tablenotes}
\end{threeparttable}
\end{table}

\casehead{Plus factors and the cost of communication}
Communication helps a cartel, but it also creates evidence. Price-fixing law in the United States and the European
Union condemns agreements, and proof of an agreement usually rests on proof of communication
(\cref{sec:12-enforcement}). Economists who assist courts therefore look for \emph{plus factors}, the body of
economic circumstantial evidence of collusion beyond the parallel movement of prices. A plus factor that delivers
a strong inference of collusion is a \emph{super-plus factor} \citep[p.~213]{MarshallMarx2012}. Marshall and Marx
organize plus factors by the collusive structure that generates them, in nine baskets of cartel conduct
(pp.~213--214). Pricing structures produce four of them: (1) prices elevated above their level without the
conspiracy and (2) industry output restricted below that level. They also produce (3) steps to reduce buyers'
resistance to price increases and (4) changes in incentives within member firms that inhibit rivalry between them. Allocation structures produce
(5) an allocation of the collusive gain among members and (6) redistributions of gains and losses that maintain
compliance, such as end-of-period purchases from one another (p.~106). Enforcement structures produce
(7) communication and monitoring, especially of production, sales and market shares, and (8) a readiness to
abandon collusion if members keep cheating. A cartel that has suppressed rivalry among its members may finally seek
further profits through (9) dominant-firm conduct. Courts often ask in addition for actions that would be against a
firm's interest if taken independently \citep[p.~110]{Kaplow2013}.

Marshall and Marx measure the evidential value of a plus factor with Bayes' rule (pp.~234--235). Let $C$ be the event
that the firms collude and $F$ the event that a plus factor is observed. Their equation (11.1) is
\begin{equation}\label{eq:12-plus-bayes}
\Pr[C \mid F] = \frac{\Pr[F \mid C]\,\Pr[C]}{\Pr[F \mid C]\,\Pr[C] + \Pr[F \mid \text{not } C]\,\Pr[\text{not } C]} .
\end{equation}
An observed conduct is a plus factor if $\Pr[C \mid F] > \Pr[C]$. Define the \emph{strength} of the plus factor and
the baseline odds against a cartel as
\begin{equation*}
S = \frac{\Pr[F \mid C]}{\Pr[F \mid \text{not } C]} ,
\qquad
O = \frac{\Pr[\text{not } C]}{\Pr[C]} .
\end{equation*}
Divide the numerator and the denominator of \eqref{eq:12-plus-bayes} by $\Pr[F \mid C]\Pr[C]$. The second term of the
denominator becomes $\frac{O}{S}$, so
\begin{equation}\label{eq:12-plus-strength}
\Pr[C \mid F] = \frac{1}{1 + \frac{O}{S}}
\end{equation}
(p.~235, note~42). The posterior rises with $S$, and the conduct is a plus factor if and only if $S > 1$. A
super-plus factor is one that is unlikely without a cartel. When $\Pr[F \mid \text{not } C]$ is close to zero, $S$ is
large and $\Pr[C \mid F]$ is close to one. Marshall and Marx give an example (p.~236, note~44). Fit an econometric
model to noncollusive benchmark data and build a 95 percent \emph{prediction band} for the actual price path. A
prediction band covers realized prices, so it adds the noise in prices to the estimation error. Here it must cover
the whole path at once with probability 0.95 when there is no cartel. Let $F$ be the event that some actual price
leaves the band, so that $\Pr[F \mid \text{not } C] = 0.05$. Marshall and Marx speak of confidence bounds. A
confidence band for the mean price would not give this probability, because realized prices scatter around the mean
even without a cartel. With
$\Pr[F \mid C] = 0.99$ and baseline odds $O = 2$, that is $\Pr[C] = \frac{1}{3}$,
the strength is $S = 19.8$ and $\Pr[C \mid F] = \frac{1}{1 + \frac{2}{19.8}} \approx 0.908$. Plus factors that are weak
alone can be strong together, and $S$ extends to groups of plus factors (p.~235).

Marshall and Marx list six super-plus factors (pp.~235--236).
\begin{enumerate}[label=(\arabic*), leftmargin=*, itemsep=2pt]
\item Transaction prices above predicted levels. A reliable econometric model accounts for all material
noncollusive effects and is estimated on benchmark data where conduct was presumably noncollusive. Its predicted
prices are inconsistent with the actual path at a specified confidence level.
\item Communication and information sharing. A firm knows the details of another firm's transactions, production,
sales or inventories, which the latter would be disadvantaged to convey unilaterally.
\item Interfirm transfers that are largely void of productive noncollusive motives.
\item Within-firm incentives. In an industry with a largely homogeneous product, several firms change the incentives
of their sales forces at the same time. The incentives shift from the pursuit of market share to the maintenance of
high prices.
\item Dominant-firm conduct by nondominant firms. A subset of firms with a large enough joint share engages jointly
in dominant-firm conduct, although no single member has the market power to act as a dominant firm.
\item Economic evidence in combination. While prices and profits are high and rising, a subset of firms restricts
production, or keeps stable market shares, customer incumbency or geographic dominance despite excess capacity.
\end{enumerate}
Cartels respond (pp.~236--237). If courts treat a conduct as a super-plus factor, colluding firms try to avoid it,
but avoidance is costly. A cartel that stops making interfirm transfers must police market shares more tightly, with
no deviations left at year end. A cartel that manufactures volatility in market shares to hide their stability can
destabilize itself, especially if it rests on a market-share agreement. The econometric screens of
\cref{ch:collusion-evidence} are attempts to produce the first super-plus factor.

\citet{Kaplow2013} questions why the law should focus on communication at all (\cref{sec:12-kaplow}). The theory of
this section shows when communication matters most. It matters when firms hold private information that the
agreement needs.

\subsection{Facilitating practices and market structure}\label{sec:12-facilitating}

A \emph{facilitating practice} is a business practice that makes collusion easier to sustain without itself being
an agreement on price. Such practices work on the two sides of the incentive constraint of \cref{sec:12-perfect}.
Some lower the one-period gain from a deviation. Others make a deviation more likely to be detected or the
punishment more severe. Market structure acts through the same channels, by way of the number of firms, the
frequency of interaction and the pattern of contact across markets.

\casehead{Information exchange and price transparency}
Exchanges of information about prices, sales or capacities raise the probability that a secret price cut is
detected. \Cref{prop:13-transparency} formalizes the effect. In the $N$-firm Bertrand supergame with Nash
reversion, let a deviation be detected with probability $\mu_d$ (written $\mu$ in \cref{ch:collusion-evidence}).
Collusion at the monopoly price is then sustainable if and only if
\begin{equation}\label{eq:12-transparency-recap}
  \delta \;\ge\; \frac{N-1}{N-1+\mu_d} .
\end{equation}
The threshold falls from 1 as $\mu_d \to 0$ to $1 - \frac{1}{N}$ at $\mu_d = 1$. With noisy signals, transparency also
raises the value of collusion far above the threshold, because the cartel needs fewer punishments on the path
(\cref{prop:12-an}). Transparency can also cut the other way, because buyers who see prices respond more strongly to
a deviator's price cut (\cref{rem:13-transparency-buyers}). Danish concrete prices rose after a competition authority began to publish
firm-level transaction prices (\cref{sec:13-concrete}). The airline tariff case concerned fares that carriers
announced in advance through a common publishing system (\cref{sec:13-atpco}). Advance notice of price changes
lets rivals answer an announced increase before it takes effect, which lowers the risk of leading a price increase
\citep[pp.~187 and~193]{HoltScheffman1987}. Under the Sugar Institute's moves system, rivals matched an announced price cut before
buyers purchased. Public cuts were therefore unprofitable, and only secret concessions could raise a refiner's
share (\cref{rem:13-moves}).
Public price information is often viewed as benign, or even desirable, because secrecy breeds suspicion.
\citet[pp.~202--203, note~56]{Kaplow2013} argues that this impression is often mistaken. Secrecy makes cheating more
attractive, while publicity eases coordination. In the \textit{Container} case, firms verified with competitors the
claims of buyers that they had been offered lower prices. Government auctions that publish bids afterward alert
bidders at once to any cheating. Public prices help buyers mainly when each buyer is small, buys rarely and finds
search costly. Information sharing among firms thus raises concerns close to those of collusion, but its effect on
the market is ambiguous \citep[p.~189, note~26]{AskerNocke2021}.

\casehead{Meeting-competition and most-favored-customer clauses}
A \emph{meeting-competition clause} is a promise by a seller to match any lower price that a customer documents
from a rival, or to release the customer from the contract. A \emph{most-favored-customer clause} is a promise to
extend to a customer any lower price that the seller grants to other customers. Both clauses look like discounts
to buyers. Both also change the payoff of a secret price cut. The meeting-competition clause turns the customers of
each firm into informants and keeps them from switching, as the next result shows.

\begin{proposition}[Meeting-competition clauses]\label{prop:12-mcc}
Consider the Bertrand supergame of \cref{prop:12-bertrand}, with $N$ firms, marginal cost $c$, monopoly profit
$\pi^m$, and collusion at $p^m$ supported by Nash reversion after a deviation. Customers are identical in size,
and on the collusive path each firm serves a fraction $\frac{1}{N}$ of them. A fraction $\chi \in [0,1]$ of each
firm's customers hold a meeting-competition clause with their current supplier. A firm can make offers to
individual customers. A covered customer who receives a lower offer shows it to its supplier, which matches it
and keeps the customer. Every deviation is observed at the end of the period. Then collusion at $p^m$ is a
subgame perfect equilibrium if and only if
\begin{equation}\label{eq:12-mcc}
  \delta \;\ge\; \underline\delta(\chi) \equiv \frac{(1-\chi)(N-1)}{1 + (1-\chi)(N-1)} .
\end{equation}
The threshold falls from $1 - \frac{1}{N}$ at $\chi = 0$ to $0$ at $\chi = 1$.
\end{proposition}

\begin{proof}
After a deviation every firm prices at $c$ forever, which is a Nash equilibrium of every continuation game. The
clauses create no gain there, because no price below $c$ is ever offered. On the path a firm earns
$\frac{\pi^m}{N}$ per period, worth $\frac{\pi^m}{N(1-\delta)}$. In a one-period deviation the best a firm can do
is to keep its own customers at $p^m$ and offer slightly less than $p^m$ to the rivals' customers. The covered
rival customers stay with their suppliers, which match. The uncovered ones, a mass $(1-\chi)\frac{N-1}{N}$,
switch. As the price cut shrinks to zero, the deviation profit approaches
\begin{equation*}
  \pi^m\left[\frac{1}{N} + (1-\chi)\frac{N-1}{N}\right] = \frac{\pi^m}{N}\bigl[1 + (1-\chi)(N-1)\bigr] .
\end{equation*}
The deviator earns zero afterwards. No deviation is profitable if and only if
\begin{equation*}
  \frac{\pi^m}{N(1-\delta)} \;\ge\; \frac{\pi^m}{N}\bigl[1 + (1-\chi)(N-1)\bigr] .
\end{equation*}
Multiply both sides by $\frac{N}{\pi^m} > 0$:
\begin{equation*}
  \frac{1}{1-\delta} \;\ge\; 1 + (1-\chi)(N-1) .
\end{equation*}
Invert both sides and subtract from 1:
\begin{equation*}
  \delta \;\ge\; 1 - \frac{1}{1 + (1-\chi)(N-1)} = \frac{(1-\chi)(N-1)}{1 + (1-\chi)(N-1)} .
\end{equation*}
The right side is increasing in $(1-\chi)(N-1)$, so it falls in $\chi$. Substituting $\chi = 0$ and $\chi = 1$
gives the limits.
\end{proof}

With $N = 4$ the threshold is $\frac{3}{4}$ without clauses, $\frac{3}{5}$ when half the customers are covered,
and $\frac{3}{8}$ when 80 percent are covered. The proposition isolates the effect on the temptation to cheat. The
clause also helps detection. A covered customer who brings a rival's offer to its supplier reveals the deviation,
so the clause raises the detection probability $\mu_d$ in \eqref{eq:12-transparency-recap}. In the limit
$\chi = 1$ no customer can be stolen. The threshold \eqref{eq:12-mcc} is then zero, so in this model the monopoly price
is an equilibrium even of the static game. \citet{Salop1986} analyzes meeting-competition and most-favored-customer
clauses as practices that credibly facilitate coordination \citep[p.~284, notes~69 and~71]{Kaplow2013}.
\citet[p.~284]{Kaplow2013} adds two qualifications. Even without a clause, a customer can convey a rival's cut to
its regular supplier, which would respond if the cut is cheating. The clause matters most when customers have
switching costs, so that they prefer to stay and use the rival's offer to extract a lower price. A clause also
requires the rival's offer to be credibly conveyed, the problem at issue in the \textit{Container} case.

The most-favored-customer clause works through commitment. Under a retroactive clause, a firm that later cuts
its price must rebate the difference to customers who bought earlier. A price cut then costs the firm a rebate on
all past sales, so the firm commits itself to soft pricing, and rivals answer with higher prices of their own
\citep{Cooper1986}. A contemporaneous clause, which extends any discount to all current customers, makes selective
discounts unprofitable. A firm that cuts its price to win one customer must cut it to all of them.
\citet[p.~284]{Kaplow2013} adds two conditions and a caveat. Other customers must have a sufficient chance of eventually
detecting the secret cut. The firms must coordinate on adopting the clauses, because a firm that abstains gains a
stronger incentive to cut, knowing that its bound rivals find punishment costly. Even industry-wide adoption has a
cost for the cartel, since a clause makes punishment more costly as well as cheating (note~70).

\citet[pp.~188--194]{HoltScheffman1987} study meeting-competition and most-favored-customer clauses together, which
they call \emph{best-price policies}, in a homogeneous-good market with advance public notice of list prices. With
both clauses in place, a rival can match a discount and keep its sales, while the discounter must extend the discount
to all its customers. A discount offered to all buyers then works like a Cournot expansion of output, so any list
price above the Cournot price invites such discounts (their Proposition~1). With identical costs and equal shares,
every list price above the competitive level and no higher than the Cournot price is a subgame-perfect equilibrium
without discounts (their Proposition~2). Competition is likely to intensify when a firm can offer discounts
selectively to some of a rival's customers (their Proposition~3). The result holds when an unmatched discount wins a
larger share of new buyers than the rival would win by matching it.

Another facilitating practice is \emph{basing-point pricing}. Every seller quotes a delivered price equal to a
common base price plus freight from an agreed basing point, whatever its own location.
Delivered prices are then identical across sellers, which removes freight absorption as a hidden discount. The
practice also standardizes the price structure. The firms then need to agree on, and adjust, only one or a few base
prices instead of a price at every location \citep[p.~282]{Kaplow2013}.

\casehead{Multimarket contact}
Firms that meet in several markets can link their behavior across them. A deviation in one market can then be
punished in all of them. \citet{BernheimWhinston1990} show that this linkage pools the incentive constraints of the
separate markets. We state the logic in a simple stationary form. Two firms meet in markets $k = 1, 2$. In market
$k$ a one-period deviation from the collusive scheme raises the deviator's profit by $g_k > 0$ in the deviation
period. The punishment that follows lowers its profit by $\ell_k > 0$ in every later period.

\begin{proposition}[Pooling incentive constraints across markets]\label{prop:12-multimarket}
If each market punishes only deviations in that market, collusion in market $k$ is sustainable if and only if
$\delta \ge \delta_k \equiv \frac{g_k}{g_k + \ell_k}$. If a deviation in either market is punished in both, collusion
in both markets is sustainable if and only if
\begin{equation}\label{eq:12-multimarket}
  \delta \;\ge\; \delta_{12} \equiv \frac{g_1 + g_2}{g_1 + g_2 + \ell_1 + \ell_2} = w\,\delta_1 + (1-w)\,\delta_2,
  \qquad w = \frac{g_1 + \ell_1}{g_1 + \ell_1 + g_2 + \ell_2} .
\end{equation}
Hence $\delta_{12}$ lies between $\delta_1$ and $\delta_2$, and it equals both when $\delta_1 = \delta_2$.
\end{proposition}

\begin{proof}
With separate punishments, a deviation in market $k$ gains $g_k$ now and loses $\ell_k$ in each later period. It
is unprofitable if and only if
\begin{equation*}
  g_k \;\le\; \sum_{t \ge 1}\delta^t\ell_k = \frac{\delta}{1-\delta}\,\ell_k .
\end{equation*}
Multiply by $1 - \delta$ and collect terms in $\delta$: $g_k \le \delta(g_k + \ell_k)$, which is
$\delta \ge \delta_k$. With joint punishments, any deviation triggers both punishments. A firm that deviates in one
market therefore loses nothing more by deviating in the other market as well, so the best deviation is a
deviation in both. It gains $g_1 + g_2$ and costs $\ell_1 + \ell_2$ per later period. The same algebra gives
$\delta \ge \delta_{12}$. For the weighted average, the definition of $\delta_k$ gives
$g_k = \delta_k(g_k + \ell_k)$. Summing over $k$,
\begin{equation*}
  g_1 + g_2 = \delta_1(g_1 + \ell_1) + \delta_2(g_2 + \ell_2) .
\end{equation*}
Dividing by $g_1 + g_2 + \ell_1 + \ell_2$ gives \eqref{eq:12-multimarket}.
\end{proof}

Suppose market 1 has $g_1 = 1$ and $\ell_1 = 3$, while market 2 has $g_2 = 3$ and $\ell_2 = 1$. Market 2 might be
one where orders are large and rare, so that one secret sale is worth a lot. Then $\delta_1 = \frac{1}{4}$,
$\delta_2 = \frac{3}{4}$ and $\delta_{12} = \frac{1}{2}$. For $\delta \in [\frac{1}{2}, \frac{3}{4})$, linking the
markets sustains collusion in market 2 that the firms could not sustain there alone. The slack in the easy market
is transferred to the hard one. When the two markets are identical, linking them changes nothing.

\citet[section~6.3.2.5]{Tirole1988} gives a version in which the markets differ in how often they meet. Two firms
collude at the monopoly price in two identical markets, each worth $\frac{\pi^m}{2}$ per firm and meeting. Market 1
meets every period and market 2 only in even periods, so market 2 alone has the discount factor $\delta^2$. If
$\delta^2 < \frac{1}{2} \le \delta$, collusion is sustainable in market 1 alone and not in market 2 alone. With
multimarket contact, the strongest temptation arises in even periods, when both markets are open, and a deviator
undercuts in both. It gains $\frac{\pi^m}{2}$ in each market and loses all future collusive profits. Full collusion
is sustainable if
\begin{equation*}
2\cdot\frac{\pi^m}{2} \;\le\; \frac{\pi^m}{2}\bigl(\delta + \delta^2 + \delta^3 + \cdots\bigr)
+ \frac{\pi^m}{2}\bigl(\delta^2 + \delta^4 + \cdots\bigr)
= \frac{\pi^m}{2}\Bigl(\frac{\delta}{1 - \delta} + \frac{\delta^2}{1 - \delta^2}\Bigr) .
\end{equation*}
Multiply by $\frac{2(1 - \delta^2)}{\pi^m}$ to get $2(1 - \delta^2) \le \delta(1 + \delta) + \delta^2$, that is,
$4\delta^2 + \delta - 2 \ge 0$, or $\delta \ge \frac{\sqrt{33} - 1}{8} \approx 0.593$ (Tirole's equation 6.9). At
$\delta = 0.6$, contact sustains full collusion in both markets, although market 2 alone supports none.
Tirole's Exercise~6.6 replaces the lower frequency by a longer information lag. Market 2 then needs $\delta^2 \ge
\frac{1}{2}$ alone, and contact lowers the requirement to $\delta \ge 0.64$ approximately.

\citet{BernheimWhinston1990} establish the irrelevance result for symmetric and independent Bertrand markets with
constant returns and perfect monitoring. The short-run gain from deviating and the long-run loss from forgone
collusion are both multiplied by the number of markets. Neither the threshold nor the value per market changes
\citep[pp.~205--206]{AskerNocke2021}. Contact matters when markets are asymmetric or firms are heterogeneous.
Asker and Nocke give an example with as many markets as firms (p.~206). The same $N$ firms compete in prices in
each of $N$ markets. Each firm has a home market in which its marginal cost is $c_L$, and its cost is $c_H > c_L$ in
the other markets. The cost gap is small, so that $p^m(c_L) > c_H$. Under \emph{efficient collusion} each firm
serves only its home market, at the price $p^m(c_L)$, and earns
$\pi_L^m = \bigl(p^m(c_L) - c_L\bigr)D\bigl(p^m(c_L)\bigr)$ per period.

If each market punishes only deviations in that market, efficient collusion fails for every discount factor. A firm
earns nothing in a foreign market on the path. By undercutting there it earns
$g = \bigl(p^m(c_L) - c_H\bigr)D\bigl(p^m(c_L)\bigr) > 0$, and Bertrand reversion in that market leaves it zero, the
same as before. With multimarket contact, a deviation anywhere triggers reversion everywhere. The best deviation
then undercuts in all $N - 1$ foreign markets at once. Reversion leaves the firm the limit-price profit
$\pi_B = (c_H - c_L)D(c_H)$ in its home market. The incentive constraint is
\begin{equation*}
\frac{\pi_L^m}{1 - \delta} \;\ge\; \pi_L^m + (N - 1)\,g + \frac{\delta\,\pi_B}{1 - \delta} .
\end{equation*}
Multiply by $1 - \delta$ and collect the terms in $\delta$:
\begin{equation*}
\delta \;\ge\; \frac{(N - 1)\,g}{(N - 1)\,g + \pi_L^m - \pi_B} .
\end{equation*}
As the cost gap shrinks, $g \to \pi_L^m$ and $\pi_B \to 0$, so the threshold tends to $\frac{N - 1}{N}$, the value that
Asker and Nocke report. Contact lets the firms divide the markets into spheres of influence. The pooled constraints
need not raise prices everywhere. \citet[ch.~6, note~22]{Tirole1988} notes that, as Bernheim and Whinston show,
contact may raise the price in one market and lower it in another.

Under imperfect monitoring, contact helps even when markets are symmetric \citep{Matsushima2001}. Return to the
noisy-signal model of \cref{asm:12-an}, replicated in $\bar k$ independent markets with independent signals. Suppose
the cartel reverts to competition everywhere, with some probability, only when the signal equals one in every market.
For $\delta$ close to one this is optimal, and the binding temptation is to deviate in a single market. Asker and
Nocke report the resulting value per firm and market (p.~206),
\begin{equation*}
v^*_{\bar k} = \frac{\bar k - (\bar k - 1 + N)\gamma\lambda}{(1 - \gamma\lambda)\,\bar k}\cdot\frac{\bar\pi}{N}
= \frac{1}{1 - \gamma\lambda}\Bigl[1 - \gamma\lambda - \frac{(N - 1)\gamma\lambda}{\bar k}\Bigr]\frac{\bar\pi}{N} .
\end{equation*}
With $\bar k = 1$ this is \eqref{eq:12-an-value}. The value rises with the number of markets and tends to the
perfect-monitoring value $\frac{\bar\pi}{N}$. Pooling the signals lets the cartel punish only after bad news in every
market, which a single deviation rarely produces. \citet{NockeStrausz2020} call this the \emph{informativeness effect}
\citep[p.~207]{AskerNocke2021}. \citet{EvansKessides1994} find higher airline fares on routes served by carriers that
meet on many other routes, and \citet{Miller2010} revisits their regression (\cref{sec:13-atpco}).

\casehead{Cross-ownership}
A firm that owns part of a rival internalizes part of the rival's profit. Consider the Bertrand duopoly with Nash
reversion, and let each firm own a share $a \in [0,1)$ of the other's profit. On the collusive path each firm's
objective is worth $(1+a)\frac{\pi^m}{2}$ per period. A deviator captures $\pi^m$ for one period while its stake
in the rival yields nothing, and both earn zero afterwards. Collusion is sustainable if and only if
\begin{equation*}
  \frac{(1+a)\,\pi^m}{2(1-\delta)} \;\ge\; \pi^m
  \quad\Longleftrightarrow\quad
  \delta \;\ge\; 1 - \frac{1+a}{2} = \frac{1-a}{2} .
\end{equation*}
Cross-ownership lowers the threshold because the deviator steals business from a firm it partly owns. Under Cournot
competition partial ownership also makes the static equilibrium less competitive, which weakens Nash reversion as a
punishment.
\citet[pp.~29--34]{Malueg1992} studies a Cournot duopoly with a one-parameter family of demand curves and finds that
the net effect depends on the curvature of demand. When demand is linear or concave, more cross-ownership lowers the
threshold. When demand is convex, it raises the threshold over a range of ownership shares. For some convex demands,
collusion is then easier with no cross-ownership than with a high but incomplete share.
The same partial internalization of rivals' profits drives unilateral effects of partial ownership
\citep{OBrienSalop2000} and the pricing model of \cref{sec:11-mc-theory}.

\casehead{Entry}
High collusive prices attract entrants, one of the four problems that \citet{McAfeeMcMillan1992} say every cartel must
solve \citep[pp.~202 and~207--209]{AskerNocke2021}. Predation and collusion can be complements, because a cartel can
deter entry by starting or threatening a price war against entrants. They can also be substitutes, because a strong
firm can earn high profits either by colluding or by driving weaker rivals out. \citet{Wiseman2017} studies a
stochastic game in which each firm's financial strength evolves with its profits and a firm in a bad enough state must
exit. Without entry, and under further conditions, an anti-folk theorem holds. For discount factors close to one,
every Nash equilibrium has a price war until only one firm remains. When entry is possible, firms may collude instead.
\citet{ScottMorton1997} studies 47 entry events on the routes of British shipping cartels around 1900. The cartel was
more likely to start a price war when the entrant was weaker, that is, smaller and less experienced. The pattern fits
the long-purse theory of predation. \citet{AskerHemphill2020} document how refiners and wholesalers in the Canadian
sugar market of the late 1880s used exclusionary vertical arrangements to keep out competition and support collusive
prices. If cartels often deter entry, the harm from collusion lasts longer than the cartel's own conduct suggests.

\casehead{Market structure}
The remaining structural determinants act through the comparative statics of \cref{sec:12-perfect}. Fewer firms,
more frequent interaction, faster detection and growing demand all lower the threshold discount factor
(\cref{prop:12-bertrand-ext}). Symmetric firms find it easier to agree, and a maverick with low cost or spare
capacity makes agreement harder (\cref{sec:12-perfect}). \Cref{sec:12-success} confronts these predictions with the
evidence on which cartels succeed.

\subsection{Harrington and Skrzypacz (2011): private monitoring, sales reports and transfers}\label{sec:12-hs}

\subsubsection{The question and the facts}

\casehead{The question}
How can a cartel sustain high prices when neither prices nor sales are publicly observed? And why did so many
prosecuted cartels adopt the same design of quotas, self-reported sales and interfirm transfers?
\citet{HarringtonSkrzypacz2011} answer both questions with one equilibrium. The second question matters for
enforcement. If the design is an equilibrium response to private information, its elements are markers that
authorities can look for. The theory of repeated games with private monitoring and communication
\citep{Compte1998,KandoriMatsushima1998} had derived folk theorems, which describe the set of equilibrium payoffs as
$\delta \to 1$. Harrington and Skrzypacz have a different aim. They construct one equilibrium, in a specific
oligopoly game, that matches the documented practices \citep[p.~2428]{HarringtonSkrzypacz2011}.

\casehead{The evidence}
The evidence is narrative. Unless noted otherwise it comes from decisions of the European Commission, compiled by
\citet{Harrington2006}. There is no econometric estimation. The model is judged by whether its equilibrium
reproduces the design of the agreements.

Lysine is an amino acid added to livestock feed. In the early 1990s the five major producers formed a global cartel
that lasted until mid-1995. Ajinomoto and Sewon proposed exclusive geographic markets. Archer Daniels Midland (ADM)
pushed for a sales-quota scheme instead, and the members accepted it. Panel A of \cref{tab:12-hs-allocations}
reports the 1992 allocation. The members also coordinated on price, as described in \cref{sec:12-anatomy}. Kanji
Mimoto of Ajinomoto prepared monthly ``scorecards''. Each company telephoned or mailed its sales volume to him, and
he distributed a spreadsheet at the quarterly meetings. Under ``guaranteed buy-ins'', a company that sold more than
its quota had to buy output from producers below quota. There were isolated reports of underreporting. Cheil told
the Commission that it had supplied ``misleading'' sales information, and Ajinomoto hid 3,500 tons from the
cartel's auditors. An internal memo read ``Hide 1,000 tons in Thailand internal business.'' The agreement was
nevertheless largely successful \citep[pp.~2425--2426]{HarringtonSkrzypacz2011}.

\begin{table}[htbp]
\centering
\caption{Sales allocations in three quota cartels}
\label{tab:12-hs-allocations}
\small
\begin{threeparttable}
\begin{tabular}{@{}lr@{\hspace{1.5em}}lr@{}}
\toprule
\multicolumn{4}{@{}l}{\textit{Panel A. Lysine market allocation, 1992 (tons)}} \\
\midrule
Company & Global & \multicolumn{1}{r}{Europe} & Global share (\%) \\
\midrule
Ajinomoto & 73,500 & \multicolumn{1}{r}{34,000} & 39.7 \\
Archer Daniels Midland & 48,000 & \multicolumn{1}{r}{5,000} & 25.9 \\
Kyowa & 37,000 & \multicolumn{1}{r}{8,000} & 20.0 \\
Sewon & 20,500 & \multicolumn{1}{r}{13,500} & 11.1 \\
Cheil & 6,000 & \multicolumn{1}{r}{5,000} & 3.2 \\
Total & 185,000 & \multicolumn{1}{r}{65,500} & 100.0 \\
\midrule
\multicolumn{2}{@{}l}{\textit{Panel B. Citric acid, 1991}} &
\multicolumn{2}{l@{}}{\textit{Panel C. International Steel Agreement, 1926}} \\
\cmidrule(r{1.5em}){1-2}\cmidrule{3-4}
Company & Share (\%) & Country & Share (\%) \\
\midrule
Haarmann and Reimer & 32.0 & Germany & 40.45 \\
Archer Daniels Midland & 26.3 & France & 31.89 \\
Jungbunzlauer & 23.0 & Belgium & 12.57 \\
Hoffmann-La Roche & 13.7 & Luxemburg & 8.55 \\
Cerestar Bioproducts & 5.0 & Saar Territory & 6.54 \\
\bottomrule
\end{tabular}
\begin{tablenotes}\footnotesize
\item Notes: Panels A, B and C reproduce Tables 1, 2 and 3 of \citet{HarringtonSkrzypacz2011}. Their sources are
European Commission decisions (Case COMP/36.545/F3 for lysine and Case COMP/E-1/36.604 for citric acid) and
\citet{Plummer1938} for steel. The column of global lysine shares and the totals are our computations. Citric acid
quotas are the average of each firm's sales in the previous three years. Steel shares are allocated market shares
by country. Rounded shares need not add to exactly 100.
\end{tablenotes}
\end{threeparttable}
\end{table}

From early 1991 to mid-1995 the five largest producers of citric acid operated a cartel. They made up about
60 percent of world production and 67 percent of production in the European Union. The quotas were market shares,
set at the average of each firm's sales over the previous three years (Panel B). Each month every company reported
its sales to an executive of Hoffmann-La Roche, who assembled the data and reported them back to the members by
telephone. Independent Swiss auditors checked the reported sales. Enforcement used a ``buy-back system''. A company
that exceeded its quota in one year had to buy product in the following year from the companies below quota. At
a meeting in November 1991 the members determined that Haarmann and Reimer had to buy 7,000 tons from ADM, and the
purchase seems to have been made. Actual production adhered very closely to the plan
\citep[p.~2426]{HarringtonSkrzypacz2011}.

The design was widespread. Cartels in carbonless paper, choline chloride, copper plumbing tubes, graphite
electrodes, plasterboard, vitamins and zinc phosphate set sales quotas and monitored them through reported sales.
For carbonless paper the Commission confirmed ex post that the sales figures exchanged at meetings were accurate.
Transfer schemes based on reported sales are documented for choline chloride, organic peroxides, sodium gluconate,
sorbates, vitamins and zinc phosphate. In choline chloride, for example, Akzo Nobel, UCB and BASF could claim
35, 28 and 15 percent of the market. The members accepted that compensation was due if these shares were
exceeded. The practice is old. The International Steel Agreement of 1926 fixed the quotas of Panel C and measured
each country's crude steel production monthly. A country whose quarterly production exceeded its quota paid \$4
per excess ton into a common fund. A country below its quota received \$2 per ton short from the fund
\citep[pp.~2427--2428]{HarringtonSkrzypacz2011}.

\casehead{Five common properties}
\citet[p.~2427]{HarringtonSkrzypacz2011} distill five properties of these agreements. First, buyers were industrial
firms that negotiated prices bilaterally, so prices and firm sales were private information. Second, the agreement
allocated sales, beyond coordinating on a common price. Third, the cartel monitored the agreement by comparing
sales with quotas. Fourth, monitoring used self-reported sales, which on the whole were not verifiable. Fifth,
the agreement was at least partly enforced by transfers. Firms that reported sales above quota made payments,
through interfirm purchases, to firms that reported sales below quota. The puzzle is then sharp. What stops a firm
from undercutting the collusive price and then underreporting its sales to avoid the transfer?

\subsubsection{The model}

We change the paper's notation in four places. We write $N$ for its $n$, $\mu(m)$ for its probability function
$\rho(m)$, $\E[m]$ for its mean $\mu$, and $\omega_0$ for its exponent $\omega$.

\begin{assumption}[Environment]\label{asm:12-hs}
\leavevmode
\begin{enumerate}[label=(\roman*)]
\item There are $N \ge 2$ risk-neutral firms with common discount factor $\delta \in (0,1)$ and common constant
marginal cost $c$. In each period of an infinite horizon, the firms simultaneously choose prices from a compact
set, and each firm supplies whatever demand it receives.
\item Market demand $m_t$ is integer valued and independently and identically distributed over time on
$\{\underline m, \underline m + 1, \dots, \bar m\}$, with $0 \le \underline m < \bar m < \infty$. Its probability
function is $\mu(m)$ and its mean is $\E[m] = \sum_{m=\underline m}^{\bar m}\mu(m)\,m$. Market demand does not depend
on prices.
\item Given $m$ and the price vector $\mathbf{p} = (p_1,\dots,p_N)$, the vector of firm sales
$\mathbf{q} = (q_1, \dots, q_N)$, with $\sum_i q_i = m$, is drawn from a probability function
$\Psi(\mathbf{q}; m, \mathbf{p})$. Each firm's sales lie in $\{0, 1, \dots, \bar q\}$ with $\bar q \le \bar m$.
\item Each firm observes only its own price and its own sales. Reports and payments, introduced below, are public.
\end{enumerate}
\end{assumption}

Market demand is perfectly inelastic, while each firm's demand responds to prices. \citet[footnote
8]{HarringtonSkrzypacz2011} read this as demand that is inelastic up to a choke price, set for example by the price
of an alternative input, and zero above it. The results then hold for prices below the choke price. Let
$\psi_i(q; m, \mathbf{p})$ be the marginal probability that firm $i$ sells $q$ units. Firm $i$'s belief about total
demand after selling $q$ units follows from Bayes' rule:
\begin{equation}\label{eq:12-hs-belief}
  \sigma_i(m \mid q, \mathbf{p}) = \frac{\mu(m)\,\psi_i(q; m, \mathbf{p})}
  {\sum_{m'=\underline m}^{\bar m}\mu(m')\,\psi_i(q; m', \mathbf{p})} .
\end{equation}
The paper imposes three conditions \citep[pp.~2429--2430]{HarringtonSkrzypacz2011}.
\begin{enumerate}[label=A\arabic*.,leftmargin=*]
\item \emph{Symmetry.} $\Psi$ is invariant to permutations of firm identities applied jointly to $\mathbf{q}$ and
$\mathbf{p}$.
\item \emph{Top demand is always possible.} $\sigma_i(\bar m \mid q, \mathbf{p}) > 0$ for all $q$ and $\mathbf{p}$.
\item \emph{Sales are good news about demand.} If $q' > q''$, then $\sigma_i(\cdot \mid q', \mathbf{p})$
first-order stochastically dominates $\sigma_i(\cdot \mid q'', \mathbf{p})$.
\end{enumerate}
A2 and A3 hold, for example, when each of the $m$ customers independently picks a firm, so that $\psi_i$ is
binomial \citep[footnote 9]{HarringtonSkrzypacz2011}.

\casehead{Timing}
Each period has four stages. In the price stage each firm chooses a nonnegative price. In the sales stage each firm
learns its own sales. In the report stage firms simultaneously send public, costless messages
$r_i \in \{\varnothing, 0, 1, \dots, \bar q\}$, where $\varnothing$ means no message. A message is read as a sales
report. In the transfer stage each firm makes a public, nonnegative payment, which is divided equally among the
other $N - 1$ firms.

The contrast with \citet{GreenPorter1984} is instructive. There firms choose quantities, the market price is a
public signal, and only quantities are hidden. Here prices and quantities are both private, and the only public
information is what firms choose to say and pay. The model also differs from the earlier work of
\citet{HarringtonSkrzypacz2007}, where sales are public, demand is fixed, and firms may charge customer-specific
prices \citep[pp.~2430--2431]{HarringtonSkrzypacz2011}.

\casehead{The static benchmark}
Firm $i$'s expected stage profit is
\begin{equation}\label{eq:12-hs-static}
  \pi_i(p_1, \dots, p_N) = \sum_{m=\underline m}^{\bar m}\mu(m)\,(p_i - c)\sum_{q=0}^{m} q\,\psi_i(q; m, \mathbf{p}) .
\end{equation}
Let $p^N(c)$ denote a symmetric Nash equilibrium price of the stage game when marginal cost is $c$.
\begin{enumerate}[label=A\arabic*.,leftmargin=*,start=4]
\item For every $c \ge 0$, the stage game has a symmetric Nash equilibrium price $p^N(c)$ that is increasing,
continuous and unbounded in $c$.
\end{enumerate}
We read ``increasing'' as strictly increasing, so the inverse function $(p^N)^{-1}$ exists, as the paper's
formulas presuppose. A sufficient condition for A4 is that $\psi_i$ depends only on pairwise price differences,
as it does when consumers have quasi-linear preferences, together with an interior equilibrium characterized by
its first-order condition. Under these conditions $p^N(c) = c + \kappa_0$ for a constant markup $\kappa_0$
\citep[p.~2431]{HarringtonSkrzypacz2011}.

\begin{example}[Logit allocation of customers]\label{ex:12-hs-logit}
Let each of the $m$ customers buy from firm $i$ with probability
$s_i(\mathbf{p}) = \frac{\exp(-\alpha p_i)}{\sum_{k=1}^{N}\exp(-\alpha p_k)}$, with $\alpha > 0$. Then firm $i$'s
expected sales are $\E[m]\,s_i(\mathbf{p})$, and its expected profit is $(p_i - c)\,\E[m]\,s_i(\mathbf{p})$. Using
$\frac{\partial s_i}{\partial p_i} = -\alpha s_i(1 - s_i)$, the first-order condition is
\begin{equation*}
  \E[m]\left[s_i - \alpha(p_i - c)\,s_i(1 - s_i)\right] = 0
  \quad\Longleftrightarrow\quad
  p_i - c = \frac{1}{\alpha(1 - s_i)} .
\end{equation*}
In a symmetric equilibrium $s_i = \frac{1}{N}$, so
\begin{equation*}
  p^N(c) = c + \kappa_0, \qquad \kappa_0 = \frac{N}{\alpha(N-1)} .
\end{equation*}
The markup does not depend on $c$, and A4 holds. The allocation is binomial for each firm, so A2 and A3 hold too.
\end{example}

\begin{remark}[Observable demand makes the problem easy]\label{rem:12-hs-observable}
Suppose total demand $m_t$ were common knowledge, either because it is fixed or because it is observed. If the
others report truthfully, a firm that misreports makes the total of reports differ from $m_t$. The other firms cannot
tell who lied. They all learn, however, that someone did. A common punishment after any total different from $m_t$
then induces truthful reports when firms are patient. Once reports are truthful, sales are effectively public, and
the scheme of \citet{HarringtonSkrzypacz2007} sustains collusion \citep[p.~2432]{HarringtonSkrzypacz2011}. The
interesting case is therefore stochastic and unobserved market demand.
\end{remark}

\subsubsection{The lysine strategy profile}

\begin{definition}[Lysine strategy profile]\label{def:12-lysine}
The profile has a collusive phase and a noncollusive phase, and play starts in the collusive phase. Let $\hat p$
be a collusive price, $z > 0$ a payment per unit of reported sales, and
$\phi: \{0, 1, 2, \dots\} \to [0,1]$ a \emph{breakdown function}. In the collusive phase:
\begin{enumerate}[label=(\alph*)]
\item Each firm sets price $\hat p$.
\item Each firm reports its sales truthfully, $r_i = q_i$.
\item If all firms reported, each firm pays $z r_i$, divided equally among the other $N - 1$ firms. If some firm
did not report, each firm pays nothing.
\item If all firms reported and each paid $z r_j$, a public randomization keeps the industry in the collusive
phase with probability $1 - \phi(\sum_{j=1}^{N} r_j)$. With probability $\phi(\sum_{j=1}^{N} r_j)$ it moves the
industry to the noncollusive phase. After a missing report or payment, the industry moves to the noncollusive
phase for sure.
\end{enumerate}
In the noncollusive phase, which is absorbing, each firm sets $p^N(c)$, makes no report and pays nothing.
\end{definition}

The profile fits the documented practices well. Firms report their sales, and transfers move from firms that
report high sales to firms that report low sales. In practice the transfers took the form of interfirm sales. As
long as the price of an interfirm sale exceeds cost, such a sale acts like a monetary transfer, although the
records do not reveal those prices. One element of the profile is missing from the case records, namely the rule
under which the cartel collapses. Even the 1926 steel agreement did not say what would happen if a country refused
to pay \$4 per excess ton. The working assumption is that egregious behavior, such as a missing payment or an
incredible report, risks the collapse of the cartel \citep[p.~2433]{HarringtonSkrzypacz2011}.

\begin{definition}[Semipublic perfect equilibrium]\label{def:12-semipublic}
A \emph{semipublic perfect equilibrium} \citep{Compte1998} is a sequential equilibrium with two properties. Each
firm's prices and payments depend only on the public history of reports and payments. Each firm's report depends
only on the public history and on its own price and sales in the current period.
\end{definition}

The solution concept keeps beliefs simple. A firm never needs to infer the private history of its rivals, because
nothing they do depends on it beyond the current period.

\subsubsection{Pricing: the transfer as a tax on business stealing}

Let $V$ and $V^N$ be a firm's expected discounted values at the start of a period in the collusive and the
noncollusive phase. Under truthful reporting, each firm's expected net transfer is zero, and each firm sells
$\frac{\E[m]}{N}$ on average by symmetry. The expected stage profit on the collusive path is therefore
\begin{equation}\label{eq:12-hs-pihat}
  \hat\pi(z) = \bigl[\hat p - c\bigr]\frac{\E[m]}{N},
\end{equation}
with $\hat p$ tied to $z$ as shown next. The static Nash profit is $\hat\pi(0) = [p^N(c) - c]\frac{\E[m]}{N}$.

\begin{proposition}[Effective marginal cost]\label{prop:12-hs-cost}
Suppose that in the collusive phase every firm reports truthfully and pays $z r_i$ after any price it sets. Then
the price-stage incentive constraint holds if and only if
\begin{equation}\label{eq:12-hs-price}
  \hat p = p^N\!\left(c + \frac{N}{N-1}\,z\right),
\end{equation}
so each firm prices as a static oligopolist with marginal cost $c + \frac{N}{N-1}z$. Under A4, for every
$\hat p > p^N(c)$ there is a unique transfer $z > 0$ that sustains $\hat p$, namely
\begin{equation}\label{eq:12-hs-z}
  z = \frac{N-1}{N}\Bigl[(p^N)^{-1}(\hat p) - c\Bigr] .
\end{equation}
\end{proposition}

\begin{proof}
Fix rivals at $\hat p$ and let firm 1 set $p_1$, with price vector $\mathbf{p}$. If firm 1 sells $q$ units in a
period with market demand $m$, it pays $zq$ and receives $\frac{z}{N-1}$ times the others' reported sales $m - q$.
With truthful reports the total of reports is $m$, whatever the split. The continuation value therefore depends on
$m$ only. Firm 1's expected payoff is \citep[equation (9)]{HarringtonSkrzypacz2011}
\begin{equation*}
  \sum_{m=\underline m}^{\bar m}\mu(m)\sum_{q=0}^{m}\psi_1(q; m, \mathbf{p})
  \left[(p_1 - c)q + z\,\frac{m - q}{N-1} - zq\right]
  + \delta\sum_{m=\underline m}^{\bar m}\mu(m)\bigl[\phi(m)V^N + (1 - \phi(m))V\bigr] .
\end{equation*}
Collect the terms in $q$ inside the bracket:
\begin{equation*}
  (p_1 - c)q - zq - \frac{z}{N-1}\,q = \left(p_1 - c - \frac{N}{N-1}\,z\right)q .
\end{equation*}
The remaining term $z\frac{m}{N-1}$ does not depend on $q$. Its expectation is $\frac{z\E[m]}{N-1}$. The payoff
becomes
\begin{equation*}
  \sum_{m}\mu(m)\sum_{q}\psi_1(q; m, \mathbf{p})\left(p_1 - c - \frac{N}{N-1}\,z\right)q
  + \frac{z\,\E[m]}{N-1}
  + \delta\sum_{m}\mu(m)\bigl[\phi(m)V^N + (1 - \phi(m))V\bigr] .
\end{equation*}
Only the first term depends on $p_1$, and it is the static profit \eqref{eq:12-hs-static} at marginal cost
$c + \frac{N}{N-1}z$. Hence $\hat p$ is optimal against $\hat p$ if and only if it is a symmetric static Nash price
at that cost, which is \eqref{eq:12-hs-price}. By A4 the map $z \mapsto p^N(c + \frac{N}{N-1}z)$ is continuous and
strictly increasing, equals $p^N(c)$ at $z = 0$, and is unbounded. Every $\hat p > p^N(c)$ is therefore reached by
exactly one $z > 0$. Inverting \eqref{eq:12-hs-price} gives $c + \frac{N}{N-1}z = (p^N)^{-1}(\hat p)$, which is
\eqref{eq:12-hs-z}.
\end{proof}

The intuition is simple \citep[p.~2435]{HarringtonSkrzypacz2011}. When a firm sells a unit that a rival would
otherwise have sold, it pays $z$ instead of receiving $\frac{z}{N-1}$. The net cost of stealing one unit is
\begin{equation*}
  z + \frac{z}{N-1} = \frac{N}{N-1}\,z .
\end{equation*}
The transfer works like a sales tax that is rebated to the industry. It removes part of the incentive to steal
business, and the collusive price is the static equilibrium price at the taxed cost. The cartel does not need
punishments to deter price cuts at all, since a price cut is never a profitable deviation at $\hat p$. The burden of enforcement shifts entirely to the reporting stage.

With a constant markup, $p^N(c) = c + \kappa_0$, the price rises one for one with the effective cost:
\begin{equation}\label{eq:12-hs-markup}
  \hat p - p^N(c) = \frac{N}{N-1}\,z .
\end{equation}
In \cref{ex:12-hs-logit} the collusive price is $\hat p = c + \frac{N}{N-1}z + \frac{N}{\alpha(N-1)}$. A cartel of
five firms that wants a price 10 cents above the static Nash price needs a transfer of 8 cents per unit of reported
sales.

\subsubsection{Reporting: why firms tell the truth}

Underreporting saves the payment $z$ on each unreported unit. The cartel must offset this temptation. Its only tool
is the breakdown function $\phi$, which makes the end of collusion more likely when total reported sales are low.
Let firm 1 have price $p_1$ and true sales $q_1$, and let the others report truthfully and pay. If firm 1 reports
$r_1$, the total of reports is $m - q_1 + r_1$. Its expected payoff is
\begin{equation}\label{eq:12-hs-W}
\begin{split}
  W(r_1; q_1, p_1) = \sum_{m=\underline m}^{\bar m}\sigma_1(m \mid q_1, \mathbf{p})\Bigl\{
  &(p_1 - c)q_1 + \frac{z}{N-1}(m - q_1) - z r_1 \\
  &+ \delta\bigl[\phi(m - q_1 + r_1)V^N + \bigl(1 - \phi(m - q_1 + r_1)\bigr)V\bigr]\Bigr\} .
\end{split}
\end{equation}
Truthful reporting requires $W(q_1; q_1, p_1) \ge W(r_1; q_1, p_1)$ for every $r_1$, $q_1$ and $p_1$. Write the
continuation term as $\delta V - \delta\phi\,(V - V^N)$ and subtract. The terms that do not depend on the report
cancel, and the condition becomes
\begin{equation}\label{eq:12-hs-truth}
  z\,(r_1 - q_1) + \delta(V - V^N)\sum_{m=\underline m}^{\bar m}\sigma_1(m \mid q_1, \mathbf{p})
  \bigl[\phi(m - q_1 + r_1) - \phi(m)\bigr] \;\ge\; 0 .
\end{equation}
The first term is the change in the payment. The second is the change in the expected loss from a breakdown. Taking
$r_1 = q_1 - 1$ gives the constraint against underreporting by one unit \citep[equation
(5)]{HarringtonSkrzypacz2011}:
\begin{equation}\label{eq:12-hs-icc}
  \delta(V - V^N)\sum_{m=\underline m}^{\bar m}\sigma_1(m \mid q_1, \mathbf{p})\bigl[\phi(m-1) - \phi(m)\bigr]
  \;\ge\; z .
\end{equation}
The right side is the payment saved. The left side is the expected increase in the probability of a breakdown,
times the value $V - V^N$ that a breakdown destroys.

\begin{lemma}[Reduction of the reporting constraints]\label{lem:12-hs-reduction}
Suppose $\phi(m-1) - \phi(m)$ is nonincreasing in $m$ for $m \le \bar m$, and A3 holds. Then truthful reporting is
preferred to every underreport, for every $q_1$ and $p_1$, if and only if \eqref{eq:12-hs-icc} holds at
$q_1 = \bar q$ for every $p_1$.
\end{lemma}

\begin{proof}[Proof sketch]
The argument follows the appendix of \citet{HarringtonSkrzypacz2011}. Fix $q_1$ and compare a report $r$ with
$r - 1$ for $r \le q_1$. Lowering the report by one more unit raises the breakdown probability by
$\phi(x - 1) - \phi(x)$ at the total $x = m - q_1 + r \le m$. Because these increments are nonincreasing, each
additional unit of underreporting raises the breakdown probability at least as much as the first unit, while it
saves the same $z$. So if underreporting by one unit does not pay, no larger underreport pays. Next compare sales
levels. A firm with lower sales puts more weight on low demand by A3, where the increments $\phi(m-1) - \phi(m)$ are
larger. The left side of \eqref{eq:12-hs-icc} is therefore largest for low $q_1$ and smallest at $q_1 = \bar q$.
The constraint at $\bar q$ implies all the others.
\end{proof}

\casehead{The breakdown function}
Harrington and Skrzypacz use the following breakdown function in the proof of their theorem
\citep[equation (3)]{HarringtonSkrzypacz2011}:
\begin{equation}\label{eq:12-hs-phi}
  \phi(x) =
  \begin{cases}
    \beta(\bar m - x)(1 - \delta) & \text{if } x \le \bar m,\\[2pt]
    (1 - \delta)^{\omega_0} & \text{if } x > \bar m.
  \end{cases}
\end{equation}
Here $x$ is the total of reports and $\omega_0 \in (0,1)$. The slope parameter $\beta > 0$ is unrelated to the taste
coefficients of the demand chapters. \Cref{fig:12-hs-breakdown} draws the function. The function falls linearly in
total reports up to the largest possible demand $\bar m$, where it reaches zero. Low totals are taken as evidence
of underreporting. Totals above $\bar m$ cannot occur under truthful reporting, so they reveal an overreport, and
the breakdown probability jumps up. As $\delta \to 1$ the jump level $(1-\delta)^{\omega_0}$ becomes large
relative to every on-path value, because $\omega_0 < 1$. The linear part satisfies the condition of
\cref{lem:12-hs-reduction}, since $\phi(m-1) - \phi(m) = \beta(1-\delta)$ is constant for $m \le \bar m$.

\begin{figure}[htbp]
\centering
\begin{tikzpicture}[x=1cm,y=1cm,>=Stealth]
  \draw[->] (0,0) -- (10.6,0) node[right] {$\sum_{j=1}^{N} r_j$};
  \draw[->] (0,0) -- (0,4.6) node[above] {\small Breakdown probability};
  \draw (0.08,4) -- (-0.08,4) node[left] {\small 1};
  \node[left] at (-0.08,0) {\small 0};
  \draw[thick] (0,3.2) -- (2.2,2.194);
  \draw[thick,dotted] (2.2,2.194) -- (3.4,1.646);
  \draw[thick] (3.4,1.646) -- (7,0);
  \draw[thick] (7.8,1.85) -- (10.1,1.85);
  \draw[densely dotted] (7,0) -- (7,1.85) -- (7.8,1.85);
  \foreach \x/\y in {0/3.2, 0.8/2.834, 1.6/2.469, 4.0/1.371, 6.2/0.366, 7.0/0}
    {\fill (\x,\y) circle (1.6pt);}
  \foreach \x in {7.8, 8.6, 9.4}
    {\fill (\x,1.85) circle (1.6pt);}
  \foreach \x in {0, 0.8, 1.6, 4.0, 6.2, 7.0, 7.8, 8.6, 9.4}
    {\draw (\x,0.08) -- (\x,-0.08);}
  \node[below] at (0,-0.08) {\small $0$};
  \node[below] at (0.8,-0.08) {\small $1$};
  \node[below] at (1.6,-0.08) {\small $2$};
  \node[below] at (2.8,-0.12) {\small $\cdots$};
  \node[below] at (4.0,-0.08) {\small $\underline m$};
  \node[below] at (6.1,-0.08) {\small $\bar m - 1$};
  \node[below] at (7.0,-0.08) {\small $\bar m$};
  \node[below] at (7.95,-0.08) {\small $\bar m + 1$};
  \node[below] at (9.4,-0.12) {\small $\cdots$};
  \node[right,font=\small] at (0.15,3.55) {$\beta\,\bar m\,(1-\delta)$};
  \node[font=\small,anchor=south west] at (4.6,1.25) {slope $-\beta(1-\delta)$};
  \node[font=\small,anchor=south] at (9.0,1.9) {$(1-\delta)^{\omega_0}$};
  \draw[decorate,decoration={brace,amplitude=5pt,mirror}] (4.0,-0.7) -- (7.0,-0.7)
    node[midway,below=6pt,font=\footnotesize,align=center] {totals possible\\under truthful reports};
\end{tikzpicture}
\caption{The breakdown function \eqref{eq:12-hs-phi}: the probability that collusion ends as a function of total
reported sales. On the equilibrium path the total equals market demand, which lies in
$\{\underline m, \dots, \bar m\}$. The probability falls by $\beta(1-\delta)$ per unit of reported sales up to
$\bar m$, and jumps to $(1-\delta)^{\omega_0}$ for totals that exceed the largest possible demand. Redrawn and
annotated from \citet{HarringtonSkrzypacz2011}, Figure 1.}
\label{fig:12-hs-breakdown}
\end{figure}

\needspace{6\baselineskip}
\casehead{The value of collusion}
\begin{lemma}[Value of the collusive phase]\label{lem:12-hs-value}
Under the lysine strategy profile with truthful reports \citep[equation (7)]{HarringtonSkrzypacz2011},
\begin{equation}\label{eq:12-hs-value}
  V - V^N = \frac{\bigl[\hat p - p^N(c)\bigr]\frac{\E[m]}{N}}{1 - \delta + \delta\sum_{m=\underline m}^{\bar m}\mu(m)\phi(m)} .
\end{equation}
With the breakdown function \eqref{eq:12-hs-phi}, the expected breakdown probability on the path is
\citep[equation (8)]{HarringtonSkrzypacz2011}
\begin{equation}\label{eq:12-hs-hazard}
  \sum_{m=\underline m}^{\bar m}\mu(m)\phi(m) = \sum_{m=\underline m}^{\bar m}\mu(m)\,\beta(\bar m - m)(1-\delta)
  = \beta\bigl(\bar m - \E[m]\bigr)(1 - \delta) .
\end{equation}
\end{lemma}

\begin{proof}
The noncollusive phase repeats the static equilibrium, so $V^N = \hat\pi(0) + \delta V^N$. In the collusive phase
the firm earns $\hat\pi(z)$ and then continues in the collusive phase unless a breakdown occurs:
\begin{equation*}
  V = \hat\pi(z) + \delta\sum_{m}\mu(m)\bigl[\phi(m)V^N + (1 - \phi(m))V\bigr] .
\end{equation*}
Subtract the equation for $V^N$:
\begin{equation*}
  V - V^N = \hat\pi(z) - \hat\pi(0) + \delta\Bigl(1 - \sum_{m}\mu(m)\phi(m)\Bigr)(V - V^N) .
\end{equation*}
Collect the terms in $V - V^N$:
\begin{equation*}
  (V - V^N)\Bigl[1 - \delta + \delta\sum_{m}\mu(m)\phi(m)\Bigr] = \hat\pi(z) - \hat\pi(0) .
\end{equation*}
By \eqref{eq:12-hs-pihat}, $\hat\pi(z) - \hat\pi(0) = [\hat p - p^N(c)]\frac{\E[m]}{N}$, which gives
\eqref{eq:12-hs-value}. For \eqref{eq:12-hs-hazard}, every on-path total satisfies $m \le \bar m$, so the linear
branch of \eqref{eq:12-hs-phi} applies, and $\sum_m \mu(m)(\bar m - m) = \bar m - \E[m]$.
\end{proof}

The value of collusion rises with the collusive price and falls with the breakdown rate. Here lies the tension
\citep[p.~2437]{HarringtonSkrzypacz2011}. A higher $z$ supports a higher $\hat p$, which raises $V - V^N$. But $z$
is also the saving from underreporting a unit, so a higher $z$ strengthens the temptation that \eqref{eq:12-hs-icc}
must contain. The constraint can hold only if the collusive gain is large relative to $z$, and only if breakdowns
are rare enough to keep $V - V^N$ large.

\subsubsection{The main result}

\begin{theorem}[Harrington and Skrzypacz, Theorem 1]\label{thm:12-hs}
Let \cref{asm:12-hs} and A1 to A4 hold. For any $\hat p > p^N(c)$ and any $\varepsilon_0 > 0$ there is
$\underline\delta < 1$ such that the following is true. If $\delta > \underline\delta$ and
\begin{equation}\label{eq:12-hs-cond}
  \frac{\E[m]}{\bar m - \E[m]} \;>\; \frac{(N-1)\bigl[(p^N)^{-1}(\hat p) - c\bigr]}{\hat p - p^N(c)},
\end{equation}
then the lysine strategy profile is a semipublic perfect equilibrium, and
$\max\{\phi(m) : \underline m \le m \le \bar m\} < \varepsilon_0$.
\end{theorem}

Condition \eqref{eq:12-hs-cond} is the paper's condition (1). The right side is positive, because
$\hat p > p^N(c)$ implies $(p^N)^{-1}(\hat p) > c$. The left side grows without bound as $\E[m]$ approaches
$\bar m$. The condition therefore says that demand must not be too volatile, in the restricted sense that average
demand is close to maximal demand. The proof combines the steps below.

\casehead{Step 1: the binding constraint}
With \eqref{eq:12-hs-phi}, the increment $\phi(m-1) - \phi(m)$ equals $\beta(1-\delta)$ for every
$m \in \{\underline m, \dots, \bar m\}$. Beliefs sum to one, so the left side of \eqref{eq:12-hs-icc} equals
$\delta(V - V^N)\beta(1-\delta)$ whatever the firm's sales and price. By \cref{lem:12-hs-reduction}, all
underreporting constraints reduce to
\begin{equation}\label{eq:12-hs-step1}
  \delta\,\beta\,(1-\delta)\,(V - V^N) \;\ge\; z .
\end{equation}

\casehead{Step 2: substitute the value of collusion}
Write $\Delta \equiv \hat p - p^N(c)$ for the price gap in this subsection only. By \cref{lem:12-hs-value}, the
denominator of \eqref{eq:12-hs-value} is $1 - \delta + \delta\beta(\bar m - \E[m])(1 - \delta)$, so
\begin{equation*}
  V - V^N = \frac{\Delta\,\frac{\E[m]}{N}}{(1-\delta)\bigl[1 + \delta\beta(\bar m - \E[m])\bigr]} .
\end{equation*}
Substitute into \eqref{eq:12-hs-step1} and cancel $1 - \delta$:
\begin{equation*}
  \frac{\delta\beta\,\Delta\,\frac{\E[m]}{N}}{1 + \delta\beta(\bar m - \E[m])} \;\ge\; z .
\end{equation*}
Multiply by the positive denominator and by $N$:
\begin{equation*}
  \delta\beta\,\Delta\,\E[m] \;\ge\; Nz + \delta\beta\,Nz\,(\bar m - \E[m]) .
\end{equation*}
Collect the terms in $\delta\beta$ \citep[equation (26)]{HarringtonSkrzypacz2011}:
\begin{equation}\label{eq:12-hs-step2}
  \delta\beta\Bigl[\Delta\,\E[m] - Nz\,(\bar m - \E[m])\Bigr] \;\ge\; Nz .
\end{equation}

\casehead{Step 3: when a slope exists}
The right side of \eqref{eq:12-hs-step2} is positive. If the bracket is zero or negative, no $\beta > 0$ satisfies
the constraint. If the bracket is positive, every slope at or above
\begin{equation}\label{eq:12-hs-betamin}
  \beta_{\min} \equiv \frac{Nz}{\delta\bigl[\Delta\,\E[m] - Nz\,(\bar m - \E[m])\bigr]}
\end{equation}
satisfies it. This is the paper's bound (4), since $Nz = (N-1)[(p^N)^{-1}(\hat p) - c]$ by \eqref{eq:12-hs-z}.
Within this family of breakdown functions, the underreporting constraints can therefore be met if and only if
\begin{equation}\label{eq:12-hs-step3}
  \frac{\Delta\,\frac{\E[m]}{N}}{z} \;>\; \bar m - \E[m],
\end{equation}
which is the paper's condition (6). The left side is the per-firm gain from collusion per unit of transfer. The
right side is the gap between maximal and average demand.

\casehead{Step 4: back to the price}
Multiply both sides of \eqref{eq:12-hs-step3} by $Nz$ and substitute $Nz = (N-1)[(p^N)^{-1}(\hat p) - c]$ on the
right:
\begin{equation*}
  \Delta\,\E[m] \;>\; (N-1)\bigl[(p^N)^{-1}(\hat p) - c\bigr](\bar m - \E[m]) .
\end{equation*}
Divide both sides by $\Delta(\bar m - \E[m]) > 0$ to obtain \eqref{eq:12-hs-cond}.

\casehead{Step 5: probabilities, overreports and payments}
Three loose ends remain. First, $\beta_{\min}$ stays bounded as $\delta \to 1$, so every value of $\phi$ on the
linear branch is at most $\beta\bar m(1-\delta)$, which goes to zero. This delivers both $\phi \le 1$ and the bound
$\varepsilon_0$ on the path. Second, an overreport raises the payment, and with positive probability it pushes the
total above $\bar m$ (by A2, $\sigma_1(\bar m \mid q_1, \mathbf{p}) > 0$). There the breakdown probability jumps
to $(1-\delta)^{\omega_0}$. As $\delta \to 1$ this jump outweighs the small reduction in $\phi$ on the linear branch,
so overreports do not pay. Third, a firm that reports and then fails to pay triggers a breakdown for sure. Reporting
zero and paying nothing is better than that, and truthful reporting is at least as good as reporting zero. So the
prescribed payment is optimal \citep[pp.~2447--2449]{HarringtonSkrzypacz2011}.

\begin{remark}[Overreports at the smallest slope]\label{rem:12-hs-overreport}
At the smallest admissible slope, the payment alone deters overreports, with or without the jump. This is our
observation, which the paper does not make. At $\beta = \beta_{\min}$, \eqref{eq:12-hs-step1} binds, so
$\delta\beta(1-\delta)(V - V^N) = z$. In state $m$, an overreport by $k$ units lowers the breakdown probability by
at most $\beta(1-\delta)\min\{k, \bar m - m\}$. It would lower it by exactly that much if $\phi$ were zero above
$\bar m$. The expected benefit is therefore at most $z\,\E_{\sigma}[\min\{k, \bar m - m\}]$, where $\E_{\sigma}$
takes expectations under the firm's beliefs, and the cost is $zk$. A2 puts positive weight on
$m = \bar m$, where the benefit is zero, so the benefit falls short of $zk$. The jump to $(1-\delta)^{\omega_0}$
is what makes overreports unprofitable for every steeper slope once $\delta$ is close to one. Footnote 17 of the
paper adds a practical reason. Invoices or customer lists can prove that a sale was made, but nothing can prove
that a sale was not made, so underreporting is the harder problem.
\end{remark}

\casehead{The cost of truthful reporting}
Normalize the value of collusion by $1 - \delta$. Step 2 gives
\begin{equation*}
  (1-\delta)(V - V^N) = \frac{\Delta\,\frac{\E[m]}{N}}{1 + \delta\beta(\bar m - \E[m])} .
\end{equation*}
At the most profitable admissible slope $\beta_{\min}$,
\begin{equation*}
  1 + \delta\beta_{\min}(\bar m - \E[m]) = 1 + \frac{Nz(\bar m - \E[m])}{\Delta\,\E[m] - Nz(\bar m - \E[m])}
  = \frac{\Delta\,\E[m]}{\Delta\,\E[m] - Nz(\bar m - \E[m])} .
\end{equation*}
Substituting \citep[p.~2438]{HarringtonSkrzypacz2011},
\begin{equation}\label{eq:12-hs-normalized}
  (1-\delta)(V - V^N) = \frac{1}{N}\Bigl[\Delta\,\E[m] - Nz\,(\bar m - \E[m])\Bigr] .
\end{equation}
The first term, $\frac{\Delta\E[m]}{N}$, is the per-period gain from collusion if breakdowns never occurred. The
second term, $z(\bar m - \E[m])$ per firm, is the expected cost of the breakdowns that keep reports truthful. The cost
does not depend on $\delta$. As firms become patient the breakdown probability goes to zero, but the breakdowns
become correspondingly more costly, and the loss stays bounded away from zero. Each member in effect pays $z$ per
period for each unit by which expected demand falls short of maximal demand. Condition \eqref{eq:12-hs-step3} is
exactly the condition that \eqref{eq:12-hs-normalized} be positive.

\begin{corollary}[Constant markup (Harrington and Skrzypacz, Corollary 2)]\label{cor:12-hs}
Suppose $p^N(c) = c + \kappa_0$. Then condition \eqref{eq:12-hs-cond} becomes
\begin{equation}\label{eq:12-hs-cond2}
  \frac{\E[m]}{\bar m - \E[m]} \;>\; N - 1
  \quad\Longleftrightarrow\quad
  \frac{\E[m]}{\bar m} \;>\; \frac{N-1}{N} .
\end{equation}
If it holds, then for any $\hat p > p^N(c)$ and $\varepsilon_0 > 0$ the conclusion of \cref{thm:12-hs} holds for
$\delta$ close to one. The scheme retains the fraction
\begin{equation}\label{eq:12-hs-retained}
  1 - (N-1)\,\frac{\bar m - \E[m]}{\E[m]}
\end{equation}
of the first-best collusive gain, whatever $\hat p$ and $\delta$. If market demand is uniform on
$\{\underline m, \dots, \bar m\}$, \eqref{eq:12-hs-cond2} holds for $N = 2$ if and only if $\underline m > 0$, and
for $N \ge 3$ if and only if
\begin{equation}\label{eq:12-hs-uniform}
  \frac{\bar m}{\underline m} \;<\; \frac{N}{N-2} .
\end{equation}
\end{corollary}

\begin{proof}
With $p^N(c) = c + \kappa_0$, the inverse is $(p^N)^{-1}(\hat p) = \hat p - \kappa_0$. Hence
$(p^N)^{-1}(\hat p) - c = \hat p - c - \kappa_0 = \hat p - p^N(c)$, and the right side of \eqref{eq:12-hs-cond}
equals $N - 1$. For the second form, $\E[m] > (N-1)(\bar m - \E[m])$ rearranges to $N\E[m] > (N-1)\bar m$. For the
retained fraction, divide \eqref{eq:12-hs-normalized} by $\frac{\Delta\E[m]}{N}$ and use
$Nz = (N-1)\Delta$ from \eqref{eq:12-hs-markup}:
\begin{equation*}
  1 - \frac{Nz(\bar m - \E[m])}{\Delta\,\E[m]} = 1 - (N-1)\frac{\bar m - \E[m]}{\E[m]} .
\end{equation*}
For the uniform case, $\E[m] = \frac{\underline m + \bar m}{2}$ and $\bar m - \E[m] = \frac{\bar m - \underline m}{2}$.
Condition \eqref{eq:12-hs-cond2} becomes
\begin{equation*}
  \frac{\underline m + \bar m}{\bar m - \underline m} > N - 1
  \quad\Longleftrightarrow\quad
  \underline m + \bar m > (N-1)\bar m - (N-1)\underline m
  \quad\Longleftrightarrow\quad
  N\underline m > (N-2)\bar m .
\end{equation*}
For $N = 2$ the last inequality reads $\underline m > 0$. For $N \ge 3$ divide by $(N-2)\underline m$, which is
positive because the inequality forces $\underline m > 0$.
\end{proof}

With a constant markup the transfer cancels from the condition \citep[footnote 19]{HarringtonSkrzypacz2011}. The
smallest slope also simplifies. Substituting $Nz = (N-1)\Delta$ into \eqref{eq:12-hs-betamin} gives
\begin{equation}\label{eq:12-hs-betamin2}
  \beta_{\min} = \frac{N-1}{\delta\bigl[\E[m] - (N-1)(\bar m - \E[m])\bigr]},
\end{equation}
which does not depend on the target price either. The required fraction $\frac{N-1}{N}$ in \eqref{eq:12-hs-cond2}
is the same number as the Bertrand threshold discount factor of \cref{prop:12-bertrand}. \Cref{tab:12-hs-volatility}
shows how fast the condition tightens with the number of firms. A duopoly tolerates any uniform demand with a
positive lower bound. With three firms the top of the support must be less than three times the bottom, and with
four firms less than twice the bottom. Harrington and Skrzypacz argue that these bounds are reasonable for the
commodities that cartels usually sell. The variation that matters is unanticipated, and the cartel can control for
anticipated cycles and trends when it sets quotas \citep[pp.~2437--2438]{HarringtonSkrzypacz2011}.

\begin{table}[htbp]
\centering
\caption{Demand volatility compatible with the lysine strategy profile under a constant markup}
\label{tab:12-hs-volatility}
\small
\begin{threeparttable}
\begin{tabular}{@{}l*{6}{@{\hspace{1.4em}}c}@{}}
\toprule
Number of firms $N$ & 2 & 3 & 4 & 5 & 6 & 7 \\
\midrule
Smallest admissible $\frac{\E[m]}{\bar m}$, equal to $\frac{N-1}{N}$ & $\frac{1}{2}$ & $\frac{2}{3}$ &
$\frac{3}{4}$ & $\frac{4}{5}$ & $\frac{5}{6}$ & $\frac{6}{7}$ \\[3pt]
Largest admissible $\frac{\bar m}{\underline m}$ (uniform demand), equal to $\frac{N}{N-2}$ & any & 3 & 2 &
$\frac{5}{3}$ & $\frac{3}{2}$ & $\frac{7}{5}$ \\
\bottomrule
\end{tabular}
\begin{tablenotes}\footnotesize
\item Notes: Our computation from \cref{cor:12-hs}. Both bounds are strict. For $N = 2$ any uniform demand with
$\underline m > 0$ qualifies. The script \texttt{hs\_example.py} in \texttt{figures/ch12/src} prints these values.
\end{tablenotes}
\end{threeparttable}
\end{table}

\begin{example}[A five-firm cartel]\label{ex:12-hs-numbers}
Let $N = 5$, as in lysine, with a constant markup and market demand uniform on $\{80, \dots, 100\}$. Then
$\E[m] = 90$ and $\bar m - \E[m] = 10$. Condition \eqref{eq:12-hs-cond2} holds, since $\frac{90}{10} = 9 > 4$.
Equivalently, $\frac{\bar m}{\underline m} = \frac{5}{4} < \frac{5}{3}$. By \eqref{eq:12-hs-retained} the scheme
retains the fraction
\begin{equation*}
  1 - 4\cdot\frac{10}{90} = \frac{5}{9} \approx 0.56
\end{equation*}
of the first-best collusive gain. By \eqref{eq:12-hs-betamin2} the smallest slope is
$\beta_{\min} = \frac{4}{50\delta} = \frac{0.08}{\delta}$. The expected breakdown probability per period follows
from \eqref{eq:12-hs-hazard}:
\begin{equation*}
  \beta_{\min}\,(\bar m - \E[m])\,(1 - \delta) = \frac{0.8\,(1-\delta)}{\delta} .
\end{equation*}
With $\delta = 0.95$ this is 0.042, so the collusive phase lasts 23.75 periods on average, and the largest
on-path breakdown probability (at $m = 80$) is 0.084. With $\delta = 0.99$ the expected breakdown probability
is 0.0081, and the collusive phase lasts 123.75 periods on average. In both cases the normalized gain is
$\frac{5}{9}$ of its first-best value. We checked all reporting constraints numerically, with a binomial
allocation of customers and $\omega_0 = \frac{1}{2}$. At $\beta_{\min}$ every constraint holds, and the
underreporting constraint binds. With a slope of $1.5\beta_{\min}$ and $\delta = 0.95$, a firm that sold nothing
would gain by reporting four units, so the overreporting constraint fails. With the same slope and
$\delta = 0.99$ every constraint holds again. The script \texttt{figures/ch12/src/hs\_example.py} prints these
numbers.
\end{example}

\subsubsection{Interpretation, robustness and limitations}

\casehead{Price wars on the equilibrium path}
Breakdowns occur with positive probability on the path, after low reported sales. Nobody cheats, and everybody
knows it. The breakdowns are nevertheless needed, since without them every firm would underreport. In this
respect the equilibrium resembles \citet{GreenPorter1984}, where wars follow low prices that nobody caused
(\cref{prop:12-gp-bertrand}). With infinite reversion a breakdown ends the cartel. With finite reversion the same
mechanism produces periodic price wars \citep[footnote 15]{HarringtonSkrzypacz2011}. The answer to the puzzle of
credible reports is now clear. An unverifiable report is credible because a low total raises the chance that the
valuable collusive phase ends. An inflated report that pushes the total beyond $\bar m$ reveals itself and raises
the breakdown risk. Below $\bar m$, an inflated report costs payments that exceed the breakdown risk it saves, at
least at the slope the cartel prefers (\cref{rem:12-hs-overreport}).

\casehead{Less frequent reports and transfers}
In practice reports were often less frequent than sales, and transfers less frequent than reports. Delayed
transfers change little, because a firm's current incentives depend only on the change in accumulated transfers
that its current actions cause. Suppose instead that firms report and make transfers every $\tau$ periods. The
theorem still holds with condition \eqref{eq:12-hs-cond} replaced by \citep[p.~2441]{HarringtonSkrzypacz2011}
\begin{equation}\label{eq:12-hs-tau}
  \frac{\E[m]}{\bar m - \E[m]} \;>\; \frac{(N-1)\bigl[(p^N)^{-1}(\hat p) - c\bigr]}{\hat p - p^N(c)}
  \cdot\frac{\delta^\tau - \delta^{\tau+1}}{\delta - \delta^{\tau+1}} .
\end{equation}
The last factor equals 1 when $\tau = 1$. Our own algebra, writing
$1 - \delta^\tau = (1-\delta)(1 + \delta + \dots + \delta^{\tau-1})$, shows that it equals
\begin{equation*}
  \frac{\delta^{\tau-1}}{1 + \delta + \dots + \delta^{\tau-1}},
\end{equation*}
which is below 1 for $\tau \ge 2$ and tends to $\frac{1}{\tau}$ as $\delta \to 1$. Less frequent reporting
therefore relaxes the volatility condition in this sense, because collusion is assured between reports. Why firms
delayed reports remains open. Larger transfers may attract the attention of antitrust authorities, and delay in
information can improve the efficiency of equilibria in related repeated games
\citep[p.~2441]{HarringtonSkrzypacz2011}.

\casehead{Limitations stated by the authors}
\citet[pp.~2438--2441]{HarringtonSkrzypacz2011} discuss four limitations. First, market demand is perfectly
inelastic, so total sales carry no information about prices. With elastic demand, total sales would be a signal of
prices, and better equilibria may exist. The authors conjecture that the lysine profile would remain an equilibrium,
with prices given by a modified static game. Inelastic demand is plausible for inputs such as lysine and citric
acid, which are a small part of the buyer's cost. It is plausible at least over the range of prices that would not
alert buyers to a cartel.
The assumption fails when the cartel is not all-inclusive. Outside supply, such as the Chinese producers of
vitamin C, makes total cartel sales sensitive to cartel prices. Second, condition \eqref{eq:12-hs-cond} may fail.
The online appendix gives two examples. In one, a duopoly with a positive lower bound on demand uses quantity quotas
of $\frac{\underline m}{2}$, compensates shortfalls and punishes only off the path. In the other, with
$m \in \{0, 1, 2\}$, condition \eqref{eq:12-hs-cond} is tight for semipublic equilibria without delay. Third,
some cartels hired auditors, which the model omits. Auditors provide only a noisy signal, because members can hide
sales, as Ajinomoto did. Fourth, the timing of reports differs from practice, as discussed above. In citric acid,
for example, Hoffmann-La Roche added its own report after seeing the others (footnote 10).

\casehead{Further critiques}
Three further points deserve mention. The equilibrium has no misreporting, while the lysine record shows hidden
sales and misleading reports. A version with noisy external signals of sales, or with imperfectly controlled
employees, would be needed to fit those episodes. The analysis is confined to symmetric firms and a symmetric
profile, while the documented quotas are far from equal (\cref{tab:12-hs-allocations}). Finally, the lysine profile
is one equilibrium. It need not be the best one, and \citet{AwayaKrishna2016} and others have since studied more
generally when cheap talk about private sales helps firms collude.

\casehead{Contrast with the other models of this chapter}
\citet{AtheyBagwellSanchirico2004} also study price competition with private information, but the private
information is about costs, and prices are public (\cref{sec:12-private-costs}). In their symmetric equilibria
without communication, rigid pricing removes the incentive to misrepresent costs at the price of productive
efficiency. When firms can report costs, the motive for communication is to allocate production to the lowest-cost
firm \citep{AtheyBagwell2001}. In \citet{HarringtonSkrzypacz2011}, by contrast, the motive for reporting sales is
monitoring \citep[p.~2431]{HarringtonSkrzypacz2011}. In \citet{GreenPorter1984} the public price is the monitoring
signal and nobody communicates. Here the signal is the total of cheap-talk reports, and transfers do most of the
work that punishments do elsewhere.

\casehead{Implications for enforcement}
For markets with industrial buyers and private prices, the theory points to three markers of collusion
\citep[p.~2442]{HarringtonSkrzypacz2011}. The first is the exchange of sales reports, in secret or through a trade
association. The second is interfirm sales or other compensation, because asymmetric punishments are what make the
agreement work. The third is periodic price wars, because the threat of a war is what keeps reports truthful.
Each marker alone is weak evidence. Together they form a pattern that independent conduct would rarely produce.

\casehead{Contribution}
The paper gives the first equilibrium account of the quota, report and transfer design that recurs in modern
international cartels. It connects the abstract theory of private monitoring and communication to documented
practice. It explains why cartels with private prices agree on sales allocations in addition to prices. It also
shows how a cartel can turn unverifiable statements into reliable monitoring at a bounded
cost. \Cref{sec:13-imperfect-evidence} returns to the same cases as evidence, and \cref{sec:13-asker} studies a
cartel whose internal transfers are recorded in detail.

\subsection{Tolerated cheating: a preview}\label{sec:12-bm-preview}

\begin{remark}[Tolerated cheating]\label{rem:12-bm}
In every model so far, equilibrium involves no deliberate cheating, and any detected deviation would be punished.
The records of prosecuted cartels show both deliberate price cuts and business stealing that went unpunished,
in lysine, nucleotides and citric acid among others (\cref{sec:13-imperfect-evidence}).
\citet{BernheimMadsen2017} explain this pattern. Each of two firms has a home market where its cost is lower, and
quoting a price in any market has a fixed cost. When firms are too impatient for perfect market division at the
home monopoly price, the optimal stationary agreement sets a price floor $p^{\mathrm{fl}}$. Each home firm mixes
over prices between the floor and its monopoly price, with atoms at both. The away firm occasionally raids at
prices above the floor and sometimes wins, and these raids go unpunished. A quote at or below the floor triggers a
short, harsh price war. The threat of tolerated raids keeps each home firm from raising its own price. Punishments
cannot do that job, because all available punishment power is already used to deter raids at or below the floor.
As the discount factor rises to the threshold for perfect collusion, the floor rises to the home monopoly price. The model answers the puzzle of tolerated cheating. Its environment, theorems and
numerical examples are developed in \cref{sec:13-bm}, and we do not repeat them here.
\end{remark}
